%/////////
\begin{frame}[c,fragile]{Colours}
%
The palette is the old \alert{AMSCesenaBleu}'s, and every colour takes the role its
counterpart has in \alert{AMSUniBo}: \textcolor{sbblue}{\textit{blue}} identifies,
\textcolor{sbbluedeep}{\textit{deep blue}} and \textcolor{sbgreymid}{\textit{silver}} structure
%
\begin{exampleblock}{The four that carry the style}
\begin{tabular}{@{}llll@{}}
\swatch{sbblue}     & \texttt{sbblue}     & \texttt{\#0073BA} & head line, author, bullets, alerts, links \\
\swatch{sbbluedeep} & \texttt{sbbluedeep} & \texttt{\#003F66} & block headers, emphasis, \textsc{toc} \\
\swatch{sbsilver}   & \texttt{sbsilver}   & \texttt{\#B3B3B3} & separation lines, and the bands and fills \\
\swatch{sbgreymid}  & \texttt{sbgreymid}  & \texttt{\#5A5A5A} & subtitle, date, footnotes \\
\end{tabular}
\end{exampleblock}
%
\begin{itemize}
	\item \texttt{sbblue} is \emph{French Blue}, the one colour \alert{AMSCesenaBleu} defined;
	      \texttt{sbbluedeep} is the shade it gave its block headers
	\item \texttt{sbsilver} and its tints are the greys of Beamer's \texttt{beaver}, which
	      \alert{AMSCesenaBleu} was built on
\end{itemize}
%
\end{frame}
